\section{Introduction}\label{sec:intro}

Fiat-backed stablecoins hold most of their reserves in short-dated US government debt. With aggregate supply of the coins we study near USD~\supplyEnd{}\,bn at the end of our sample, their issuers have become a visible class of buyer at the very front of the Treasury curve. The Bank for International Settlements (BIS) notes that their Treasury holdings are now ``on par with large jurisdictions and government money market funds'' and that a USD~3.5\,bn increase in stablecoin market capitalisation ``can already depress Treasury yields by around 2.5--5 basis points''~\cite{bis2025aer}. That estimate comes from Ahmed and Aldasoro~\cite{ahmed2025}, who document the flow--yield association with local projections and identify it with an instrument.

On 18~July 2025 the United States enacted the GENIUS Act~\cite{genius2025}. It turns the composition of stablecoin reserves from an issuer choice into a legal requirement, and among Treasury securities it admits only those maturing within 93~days. If stablecoin demand moves bill yields, codifying that demand could plausibly strengthen the channel, by making reserve purchases more mechanical, or concentrate it at the shortest maturities.

This paper asks whether the data so far show either. The question is narrow, and so is the design: a replication-in-spirit of the controls-only local projections of~\cite{ahmed2025}, extended with an interaction for the post-enactment regime and estimated on open data. To limit researcher degrees of freedom we committed a pre-analysis plan (PAP) with hypotheses, specification, robustness list and decision rules before downloading any yield data. A null result was declared publishable in advance.

\paragraph*{Hypotheses} Writing $\beta$ for the response of the yield to an inflow (negative if inflows lower yields), the PAP states:
\Hyp{H1}: at horizon $h=10$, $|\beta^{post}|>|\beta^{pre}|$ for the 3-month bill;
\Hyp{H2}: after enactment the response is concentrated at the 4-week and 3-month tenors, and is flatter across tenors before;
\Hyp{H3}: outflows move yields more than inflows of the same size, in both regimes.

\paragraph*{Findings} None of the three hypotheses is supported after enactment. The post-enactment response of the 3-month yield is about half the pre-enactment one, not larger, and the difference is imprecise. The short-end concentration of~H2 is already present before enactment. The outflow asymmetry of~H3 appears before enactment, where it is plausibly driven by rate-motivated redemptions during the 2022--23 tightening cycle, but not after. A pre-registered placebo break shows that the flow--yield association shifts between sub-periods even without any legislative event.

\paragraph*{Contributions} (i)~A legally dated regime break, rather than a market-size proxy, for the stablecoin--bill channel; (ii)~a term-structure test motivated by the 93-day rule; (iii)~a post-enactment sample extending seven months beyond that of~\cite{ahmed2025}; and (iv)~a fully reproducible, pre-registered pipeline in which every number in the text is generated from versioned results files.
